% Contenido del póster: Eid, Yang y Duygun (2024), The British Accounting Review. DOI 10.1016/j.bar.2024.101448
% Elaborado con apoyo de IA a partir del resumen oficial del artículo (OpenAlex), 2026-11-12.
\def\DOI{10.1016/j.bar.2024.101448}
\def\TituloEs{Cerrar la brecha de crédito: la influencia de la estructura bancaria regional en la expansión de los préstamos entre pares}
\def\TituloOriginal{Bridging the credit gap: The influence of regional bank structure on the expansion of peer-to-peer lending}
\def\Autores{Eid, Yang y Duygun}
\def\Anio{2024}
\def\Revista{The British Accounting Review}

\def\Problema{Las plataformas de préstamos entre pares crecen de forma desigual entre regiones. Una hipótesis es que llenan los vacíos que deja la banca tradicional donde su presencia es débil.}
\def\Objetivo{Investigar en qué medida la estructura regional del mercado de crédito, definida por la presencia y la capacidad de préstamo de los bancos tradicionales, explica el crecimiento de los préstamos entre pares.}
\def\Metodo{Datos de préstamos entre pares por región y un enfoque de variables instrumentales para aislar el efecto de la estructura bancaria.}
\def\Resultados{\begin{itemize}
  \item Los préstamos entre pares crecen más en las zonas desatendidas por la banca tradicional.
  \item El efecto es mayor donde los bancos pequeños tienen menos alcance.
  \item El aumento de los préstamos entre pares se asocia con un menor riesgo de impago de los prestatarios.
  \item La expansión de las plataformas mejora el bienestar financiero de los prestatarios, porque eleva su puntaje crediticio.
\end{itemize}}
\def\Conclusion{Las plataformas de préstamos entre pares funcionan como complemento de la banca: crecen donde esta atiende menos y benefician a los prestatarios.}
\def\Relevancia{La SMV publica la actividad del financiamiento participativo por plataforma y mes, pero no por región ni por receptor, de modo que en el Perú no puede verificarse si las plataformas llegan a las zonas menos atendidas por la banca. Publicar esa información permitiría replicar este análisis.}
\def\PalabrasClave{préstamos entre pares, estructura bancaria regional, brecha de crédito, variables instrumentales, inclusión}
% --- Ficha del estudio (herramientas/posteres/ficha_articulos.py) ---
\def\Ambito{Préstamos entre pares en distintas regiones}
\def\Enfoque{Empírico: variables instrumentales}
\def\FiguraPeru{1}
% --- fin de la ficha ---
